% FROZEN — do not keep editing. Living source is harsh/resume/HarshDesai_Resume.tex
% This file produced every PDF in this folder (Detectors, Instrumentation, UW-IDP v1–v6).
% Last compile was:
%   node generate-latex.mjs harsh/resume/archive/detectors-july2026/HarshDesai_July2026_Resume_Detectors.tex <new-pdf-name>
% Focus: UW IceCube / IDP Instrumentation Engineer–Mechanical Engineer (mech-first + instrumentation)
\documentclass[10pt,a4paper]{article}

\usepackage[a4paper,top=0.30in,bottom=0.28in,left=0.40in,right=0.40in]{geometry}
\usepackage{fontspec}
\setmainfont{Tinos}
\usepackage{xcolor}
\usepackage{enumitem}
\usepackage{tabularx}
\usepackage{array}
\usepackage{hyperref}
\usepackage{titlesec}
\usepackage{microtype}
\usepackage{ragged2e}

\definecolor{nameblue}{HTML}{214F78}
\definecolor{linkblue}{HTML}{0000CC}

\hypersetup{
  colorlinks=true,
  urlcolor=linkblue,
  linkcolor=linkblue,
  pdfauthor={Harsh Desai},
  pdftitle={Harsh Desai Resume --- UW IDP Mechanical Instrumentation}
}

\pagestyle{empty}
\setlength{\parindent}{0pt}
\setlength{\tabcolsep}{0pt}
\setlength{\emergencystretch}{1em}
\fontsize{8.30pt}{9.00pt}\selectfont

\titleformat{\section}
  {\fontsize{10.2pt}{10.7pt}\selectfont\bfseries\uppercase}
  {}{0pt}{}
  [\vspace{-2.2pt}\titlerule]
\titlespacing*{\section}{0pt}{5.8pt}{1.8pt}

\setlist[itemize]{
  leftmargin=13.5pt,
  label=\raisebox{0.18ex}{\tiny$\bullet$},
  itemsep=0.22pt,
  parsep=0pt,
  topsep=0.30pt,
  partopsep=0pt
}

\newcommand{\resumeSubheading}[4]{%
  \par\noindent\begin{tabular*}{\textwidth}{@{\extracolsep{\fill}}l r}
    \textbf{#1} & \textcolor{nameblue}{\textbf{#2}} \\
    #3 & #4
  \end{tabular*}\par\vspace{-1.6pt}%
}

\newcommand{\resumeProjectHeading}[2]{%
  \par\noindent\begin{tabular*}{\textwidth}{@{\extracolsep{\fill}}l r}
    \textbf{#1} & #2
  \end{tabular*}\par\vspace{-1.6pt}%
}

\newcommand{\resumeItemsStart}{\begin{itemize}}
\newcommand{\resumeItemsEnd}{\end{itemize}\vspace{-1.1pt}}
\newcommand{\resumeItem}[1]{\item #1}

\begin{document}

% Header
\begin{tabular*}{\textwidth}{@{\extracolsep{\fill}}l r}
  {\fontsize{18pt}{18pt}\selectfont\bfseries\textcolor{nameblue}{Harsh Desai}} &
  \fontsize{8.05pt}{8.5pt}\selectfont
  Ann Arbor, MI, USA $\mid$
  \href{mailto:harshdes@umich.edu}{harshdes@umich.edu} $\mid$
  +1-734-548-1080 $\mid$
  \href{https://www.linkedin.com/in/harshddes/}{LinkedIn} $\mid$
  \href{https://harshddes.github.io/}{harshddes.github.io}
\end{tabular*}
\vspace{-3.5pt}

{\itshape Mechanical engineer with hands-on experience in scientific instrumentation, electromechanical test stands, vacuum/HV systems, plasma diagnostics, DAQ automation, and mechanical design; experienced in taking experimental hardware from setup and integration through testing and troubleshooting.}
\vspace{0.5pt}

\section{Education}
\resumeSubheading
  {University of Michigan}{Ann Arbor, MI, USA}
  {Master of Engineering, Space Systems Engineering $\mid$ GPA: 3.79/4.0}{Aug 2024--Dec 2025}
{\fontsize{7.7pt}{8.2pt}\selectfont Coursework: Plasma \& Fields Instrumentation; Plasma Measurement Techniques; Control for Aerospace Vehicles; Space Systems Design \& Mgmt.}\par\vspace{2.0pt}
\resumeSubheading
  {Vellore Institute of Technology (VIT)}{Vellore, India}
  {Bachelor of Technology, Mechanical Engineering $\mid$ GPA: 7.78/10}{Jul 2019--Jul 2023}

\section{Technical Skills}
\textbf{Mechanical Design \& Analysis:} SolidWorks, Fusion 360, Autodesk Inventor, AutoCAD, ANSYS Mechanical/Fluent\\[-0.2pt]
\textbf{Instrumentation \& Diagnostics:} Langmuir probe, emissive probe, RPA, CEM/ESA calibration, SIMION, SRIM, Amptek DPPMCA\\[-0.2pt]
\textbf{DAQ \& Controls:} Python, LabVIEW (CLAD prep), PyVISA/PyMeasure, Keithley DAQ/SMU, TDK Lambda PSUs, GPIB, serial, HV switching\\[-0.2pt]
\textbf{Digital Readout:} ADC/FPGA workflow, Zmod ADC 1410, Zynq-7000, AMD Vivado, MATLAB HDL Coder\\[-0.2pt]
\textbf{Lab \& Documentation:} Vacuum chamber operation (CTI Hi-Vac / roughing to $\sim$1~mTorr), 1.3~kV test stands; Microsoft Word, Excel, PowerPoint; technical test plans and reports

\section{Engineering \& Instrumentation Experience}

\resumeSubheading
  {Climate and Space Sciences and Engineering, University of Michigan}{Ann Arbor, MI, USA}
  {Research Assistant, LVACCS Instrumentation \& Test-Stand Characterization}{Feb 2026--Present}
\resumeItemsStart
  \resumeItem{Built and operated an instrumented 1.3~kV HV test workflow with discharge-box FMEA, Python/PyVISA ignition control, PSU logging, and DAQ synchronization; cut manual steps 15$\rightarrow$5 and achieved 98.6\% synchronized data capture}
  \resumeItem{Developed a Python/PyVISA + Tkinter operator GUI to automate ignition sequences, mass-flow control, and remote PSU coordination for high-throughput test-stand operation}
  \resumeItem{Operated Langmuir, emissive, and retarding-potential probes during vacuum-chamber experiments, performing bias sweeps and interpreting I--V / plasma-potential-related measurements across regolith-lofting test states}
  \resumeItem{Configured Keithley DAQ and TDK Lambda HV sequencing for hollow-cathode discharge testing (1.3~A constant-current, 1300~V headroom, xenon ignition near 500~V / 50~SCCM) on a CTI Hi-Vac / roughing-pumped chamber}
\resumeItemsEnd

\resumeSubheading
  {Solar and Heliospheric Research Group, University of Michigan}{Ann Arbor, MI, USA}
  {Student Research Assistant, SSD Instrumentation Readout Chain}{Jan 2025--Dec 2025}
\resumeItemsStart
  \resumeItem{Developed and verified an ADC/FPGA readout architecture for a solid-state detector, sizing digitization requirements and evaluating a 14-bit, 125~MSPS Zmod ADC / Zynq-7000 path using MATLAB and Vivado}
\resumeItemsEnd

\resumeSubheading
  {Space Physics Research Lab, University of Michigan}{Ann Arbor, MI, USA}
  {Summer Intern, TestBedz Spacecraft Qualification Platform}{May 2025--Jul 2025}
\resumeItemsStart
  \resumeItem{Built a full-stack requirement-flowdown platform for TVAC, vibration, and EMI test submissions, matching spacecraft hardware profiles to facility instrumentation limits across 6 test-type workflows}
  \resumeItem{Encoded hardware metadata, facility operating envelopes, and acceptance criteria into machine-readable compatibility checks and traceable test plans for repeatable qualification workflows}
  \resumeItem{Surfaced incompatible TVAC/vibration requests before scheduling by translating facility limit envelopes into automated submission checks}
\resumeItemsEnd

\resumeSubheading
  {L3Harris and University of Michigan}{Ann Arbor, MI, USA}
  {Communications Lead, Uranian Orbiter and Probe Mission Study}{Sep 2024--May 2025}
\resumeItemsStart
  \resumeItem{Led Uranian mission communications architecture trades and DSN/fault-tolerance requirement flowdown, managing 15 traceable requirements and supporting a 16\% (\$450M) cost reduction through an orbiter-only configuration trade}
\resumeItemsEnd

\section{Mechanical Design \& Engineering Projects}
\resumeProjectHeading{CANSAT 2022 (NASA \& AAS) $\mid$ Ranked 7th worldwide out of 42 teams}{Aug 2021--Jul 2022}
\resumeItemsStart
  \resumeItem{Designed and built a 10~m unidirectional tether-deployment mechanism using a DC motor and custom worm-gear spool for non-backdrivable load transfer and controlled descent of a 2~mm braided-nylon payload line}
  \resumeItem{Integrated payload-release and camera-stabilization hardware---elastic-lid release with torsion-spring latch plus a dual-servo 2-axis gimbal---and iterated mechanical fit, assembly, and field checkout for fixed 45$^{\circ}$ imaging}
\resumeItemsEnd

\resumeProjectHeading{Space Instrumentation Calibration \& Ion-Optics Series}{Jul 2025--Dec 2025}
{\itshape Graduate Course: SPACE 571 $\mid$ Supervisor: Prof. Stefano Livi}\vspace{-1.3pt}
\resumeItemsStart
  \resumeItem{Calibrated Channel Electron Multiplier (CEM) response and the charge-injection $\rightarrow$ CSA/shaper $\rightarrow$ MCA signal chain, extracting gain, centroid, FWHM, and bias-dependent count response}
  \resumeItem{Calibrated Cylindrical Electrostatic Analyzer (ESA) performance from lab energy-angle scans and SIMION/SRIM ion-optics models, quantifying transmission, angular acceptance, and energy resolution}
  \resumeItem{Modeled SIMION dual-Einzel beam expansion and Bessel-box filter response, and used SRIM TRANSMIT to estimate carbon-foil energy loss / straggling against TOF-derived exit energies}
\resumeItemsEnd

\resumeProjectHeading{Spaceport America Cup / IREC $\mid$ Placed 23rd globally and 5th in Asia-Pacific}{Apr 2020--Jul 2021}
\resumeItemsStart
  \resumeItem{Launched a Mach~0.9 solid-motor sounding rocket to a 10,000~ft target altitude with the competition team}
  \resumeItem{Developed an SRAD trajectory simulator using CFD-backed aerodynamic data for flight prediction and design iteration}
\resumeItemsEnd

\resumeProjectHeading{CANSAT 2021 (NASA \& AAS) $\mid$ Placed 13th globally and 7th in Asia-Pacific}{Aug 2020--Jul 2021}
\resumeItemsStart
  \resumeItem{Built MATLAB telemetry / ground-control software for dual maple-seed payload deployment monitoring, and modeled mono-wing descent with Blade Element Theory plus transient CFD}
\resumeItemsEnd

\end{document}
